% Используется также в РЗТ, ДЗТ

\color{unidarkgreen}{\subsection{Защита от неполнофазного режима}}
\color{black}

%%%%%%%%%%%%%%%%%%%%%%%%%%%%%%%%%%%%%%%%%%%%%%%%%%%%%%%%%%%%%%%%%%%%%%  ОБЩИЕ СВЕДЕНИЯ %%%%%%%%%%%%%%%%%%%%%%%%%%%%%%%%%%%%%%%%%%%%%%%%%%
\begin{enumerate}

\item
Для контроля отключения всех фаз выключателя и ликвидации неполнофазного режима на защищаемой линии предусмотрена функция защиты от неполнофазного режима (ЗНР), которая вводится в работу программным ключом \textbf{«ЗНР»}.

\item
Пуск ЗНР происходит при выполнении следующих условий:
\begin{itemize}
	\item наличие сигнала срабатывания ЗНФ В1 или ЗНФ В2;
	\item срабатывание измерительного органа, действующего по расчетному току нулевой последовательности или по отношению токов обратной и прямой последовательностей -- выбор осуществляется программным переключателем \textbf{«ИО ЗНР»};
	\item отключенное положение смежного выключателя (для схем с двумя выключателями).
\end{itemize}

\item
Наличие второго выключателя определяется программным переключателем \textbf{«Кол-во выкл-ей»}. Для полуторной схемы также реализован контроль отключенного положения выключателя В3.

\item
Порог срабатывания измерительного органа по расчетному току нулевой последовательности задается уставкой \textbf{«3I0ср ЗНР»}, по отношению токов обратной и прямой последовательностей – \textbf{«I2/I1 ЗНР»}.

\item
Выдержка времени на срабатывание ЗНР регулируется уставкой \textbf{«Тср ЗНР»}. 

\item
ЗНР может быть оперативно выведена из работы от внешнего сигнала \textit{«ЗНР: Вывод»}, назначенного на соответствующий дискретный вход или функциональную клавишу.

\item
Функциональная схема логики ЗНР приведена на рисунке \ref{pict_znr:ZNR}. Параметры для настройки ЗНР представлены в таблице \ref{app_set:znr}.

\medskip
\begin{figure}[H]
\centering
\includegraphics[width=1\textwidth,height=1\textheight,keepaspectratio]{\fbpath/03.01 ЛВ/561.0221 ЗНР/img/ZNR.pdf}
\captionof{figure}{Функциональная схема логики ЗНР}\label{pict_znr:ZNR}
\end{figure}


\end{enumerate}